\documentclass[11pt,a4paper]{article}
\usepackage[margin=1in]{geometry}
\usepackage[T1]{fontenc}
\usepackage[utf8]{inputenc}
\usepackage[english]{babel}
\usepackage{amsmath,amssymb,amsthm}
\usepackage{microtype}
\usepackage{array,tabularx,booktabs}
\usepackage[hidelinks]{hyperref}
\hypersetup{pdfauthor={Nikolai Dereviankin},
 pdftitle={Reading a Past Label versus Updating a Surface Record}}
\newtheorem{proposition}{Proposition}[section]
\newtheorem{lemma}[proposition]{Lemma}
\theoremstyle{remark}
\newtheorem{remark}[proposition]{Remark}
\newcommand{\norm}[1]{\left\lVert#1\right\rVert}
\newcommand{\R}{\mathbb R}
\setlength{\emergencystretch}{2em}
\title{Reading a Past Label versus Updating a Surface Record:\\
An Explicit Bulk--Boundary Example}
\author{Nikolai Dereviankin}
\date{16 October 2026\\[0.35em]\small Version 0.2}
\begin{document}
\maketitle
\begin{abstract}
A surface observable may distinguish a selected past label without
containing enough information to predict its response to another input.
We exhibit this distinction in a massless real wave field on a half-line
coupled to a dynamical boundary coordinate. An exact reduction to a
forced linear oscillator gives two ideal equal-energy pulse-order
preparations whose scalar surface readouts remain separated on a specified
finite interval under bounded scalar observation error. On a broader reachable family,
the same scalar fails to determine its next value under a common pulse:
a missing field--rate orientation enters the update. A two-component
record relative to a supplied phase reference admits a closed update
when the next input is specified. We state the different preparation
families and access assumptions explicitly. Motivated by a broader
question in Field 01, this is a self-contained methodological example
of readable content versus predictive sufficiency, not a test of the
cosmological proposal, a new principle of state reduction, or a mechanism
of autonomous, persistent physical memory.
\end{abstract}

\section{Model, access and the question}
\label{sec:model}

\subsection{Two tasks for a record}

Consider two questions about an observable left after an incident pulse.
Can it identify which member of a declared preparation family was used?
Can it also determine its own value after a further specified pulse?
The first question concerns a chosen past label; the second concerns
closure under subsequent dynamics. They have different quantifiers.
A successful decoder on two preparations need not provide a state
representation for all possible histories.

We study these questions in a fixed linear bulk--boundary model.
The positive result is an equal-energy order code. The negative result
uses a larger family to expose information lost by a scalar observation.
A sufficient enlarged record then restores controlled closure. These
are three statements about explicitly stated families, not contradictory
claims about the same two states. The response formulas, reachability
construction and decoding margin are derived analytically below.

The question is motivated by Field 01's broader recursive-recording
programme. That programme supplies the motivation, not the equations:
the action, preparation family and observables below are specified and
analysed on their own. A successful finite-time decoder is therefore a
result about this boundary-oscillator model, not evidence for a
cosmological memory law or for the validity of Model 01.

The underlying ingredients have established precedents. A wave field
coupled to a boundary oscillator is a known model class\cite{george}.
Linear propagation with prescribed inputs is standard state-space
dynamics\cite{frazzoli}. Exact lumpability formalizes compatible state
reduction\cite{horstmeyer}, while causal-state constructions study
predictive sufficiency and its input--output version\cite{shalizi,barnett}.
The present contribution is an explicit comparison of access rules in
one solvable preparation protocol; no general novelty claim is attached
to those ingredients. Section~\ref{sec:discussion} makes the comparisons
and the limits precise.

\subsection{A real bulk field and a dynamical boundary}

Let $u(t,x)$ be real on $x\le0$, and let $q(t)$ be an independent real
boundary coordinate. The propagation speed is one. For fixed positive
parameters $K,M,g$, take the action
\begin{equation}
 S=\int dt\left\{
 \frac K2\int_{-\infty}^0(u_t^2-u_x^2)\,dx
 +\frac M2\dot q^2-\frac g2(q-u(t,0))^2\right\}.
 \label{eq:action}
\end{equation}
There is no bulk mass or separate surface potential. This is a specified
model, not an approximation obtained by dropping terms during the
calculation. One may also regard it as the tangentially constant sector
of a flat half-cylinder with finite transverse area. In that reading,
all displayed energies are per unit transverse area.

Write $v(t)=u(t,0)$ and take the outward normal at the endpoint to point
in the positive $x$ direction. Variations vanish at the temporal endpoints
and are compactly supported in space, but their boundary values are free.
Integration by parts gives the boundary coefficients
$-K u_x-g(v-q)$ for $\delta v$ and $-M\ddot q-g(q-v)$ for $\delta q$.
Thus the equations are
\begin{align}
 u_{tt}-u_{xx}&=0, &&x<0,\\
 K u_x(t,0)+g(v-q)&=0,\\
 M\ddot q+g(q-v)&=0.
 \label{eq:system}
\end{align}
The endpoint is fixed and timelike; no moving geometry is being inferred.
The boundary variable is not identified with the trace of the bulk field.

The nonnegative energy is
\begin{equation}
 E(t)=\frac K2\int_{-\infty}^0(u_t^2+u_x^2)\,dx
       +\frac M2\dot q^2+\frac g2(q-v)^2.
 \label{eq:energy}
\end{equation}
For smooth fields with no flux at spatial infinity at finite time,
\begin{align*}
 \dot E_{\rm bulk}&=K v_t u_x(t,0),\\
 \dot E_\Sigma&=M\dot q\ddot q+g(q-v)(\dot q-\dot v)
              =g(v-q)\dot v.
\end{align*}
The Robin equation cancels these two terms. The trace $v$ enters the
surface interaction energy; later we will use a surface-only readout
which is not this energy.

\subsection{Causal pulse preparations and their exact evolution}

An admissible incident profile $F$ is a real $C^3$ function of compact
support in $[0,T_F]$, extended by zero, with its derivatives through order
three vanishing at support endpoints. Finite sums of separated pulses
with this property are admitted. Prepare
\begin{equation}
 u(0,x)=F(-x),\qquad u_t(0,x)=F'(-x),\qquad q(0)=\dot q(0)=0.
 \label{eq:initial}
\end{equation}
The compatibility conditions at the boundary hold at $t=0$. All future
incident pulses in $F$ are already encoded in these initial bulk data;
there is no later boundary work used to choose their order.

Define
\begin{equation}
 \gamma=\frac gK,\qquad \omega_*^2=\frac gM,
 \qquad \Omega^2=\omega_*^2-\frac{\gamma^2}{4}.
 \label{eq:parameters}
\end{equation}
From here on assume the underdamped condition $\Omega^2>0$. The causal
kernel of $\partial_t^2+\gamma\partial_t+\omega_*^2$ is
\begin{equation}
 G(t)=e^{-\gamma t/2}\frac{\sin(\Omega t)}{\Omega}\quad(t\ge0),
 \qquad G(t)=0\quad(t<0).
 \label{eq:kernel}
\end{equation}

\begin{proposition}[Exact pulse response]
\label{prop:response}
For the data \eqref{eq:initial}, the unique classical causal solution is
\begin{align}
 q(t)&=2\omega_*^2\int_0^tG(t-s)F(s)\,ds,\\
 h(t)&=F(t)-\frac MK\dot q(t),\\
 u(t,x)&=F(t-x)+h(t+x),
 \label{eq:solution}
\end{align}
where $h$ is zero at negative arguments. In particular,
\begin{equation}
 \ddot q+\gamma\dot q+\omega_*^2q=2\omega_*^2 F.
 \label{eq:oscillator}
\end{equation}
Uniqueness is within the smooth finite-energy class with the prescribed
initial data and no additional incoming excitation.
\end{proposition}
\begin{proof}
The characteristic decomposition solves the bulk wave equation. At the
boundary $v=F+h$ and $u_x=-F'+h'$, so the Robin equation is
$\dot v=2F'-\gamma(v-q)$. The surface equation gives
$\ddot q=\omega_*^2(v-q)$. Differentiating and eliminating $v$ yields
\[
 \dddot q+\gamma\ddot q+\omega_*^2\dot q=2\omega_*^2F'.
\]
Initially $v=q=\ddot q=F=0$, so its integration constant is zero.
This proves \eqref{eq:oscillator}. The kernel has $G(0)=0$, $G'(0+)=1$,
and solves the homogeneous equation at positive times, which proves
the convolution formula. The oscillator equation then gives
$v=q+\ddot q/\omega_*^2=2F-(M/K)\dot q$, hence the stated $h$.
These identities verify both boundary equations. Zero extension and
endpoint regularity give the required initial data and classical matching.
The incoming characteristic is fixed by \eqref{eq:initial}; the ordinary
differential equation with zero initial data fixes $q$, and then $h$.
This also proves uniqueness in the declared class.
\end{proof}

The reduction is exact for the admitted data. It is not an instantaneous
kick approximation or a phenomenological insertion of damping. The
half-line permits an outgoing wave to leave without returning.

\begin{proposition}[Conserved total energy and a decaying boundary response]
\label{prop:energy}
For \eqref{eq:solution},
\begin{align}
 E_{\rm bulk}(t)&=K\int_t^\infty F'(s)^2\,ds
                   +K\int_0^t h'(s)^2\,ds,\\
 E_\Sigma(t)&=\frac M2\dot q^2+\frac{M^2}{2g}\ddot q^2,\\
 \dot E_\Sigma&=K(F'^2-h'^2),\\
 E(t)&=K\int_0^{T_F}F'(s)^2\,ds.
 \label{eq:balance}
\end{align}
After the last incident pulse, $q$, $\dot q$, $\ddot q$ and $E_\Sigma$
tend to zero as $t\to\infty$.
\end{proposition}
\begin{proof}
Since $u_t=F'+h'$ and $u_x=-F'+h'$, their squared sum is
$2(F'^2+h'^2)$. Changing variables in the half-line integrals gives
the first identity. Substituting $q-v=-\ddot q/\omega_*^2$ gives the
second. Differentiation, using the derivative of
\eqref{eq:oscillator}, gives
\begin{align*}
 \dot E_\Sigma&=2M F'\ddot q-\frac{M^2}{K}\ddot q^2\\
              &=K\bigl[F'^2-(F'-(M/K)\ddot q)^2\bigr].
\end{align*}
This is the third identity, which cancels the derivative of the bulk
energy. Evaluation at zero proves the last identity. After $T_F$ the
oscillator has roots $-\gamma/2\pm i\Omega$, so the coordinate and its
first two derivatives decay. The expression for $E_\Sigma$ finishes the proof.
\end{proof}

Thus surface decay and conservation of the full energy coexist. The
energy has entered the outgoing bulk wave rather than disappeared.
The conclusion depends on the declared geometry and preparation; a
reflecting far boundary would supply a returning input.

\subsection{Reading and updating as factorization problems}

Let $X$ be a declared set of states, $m:X\to\mathcal Z$ a selected past
label, and $\mathsf b:X\to\mathcal Y$ an accessible record. A decoder
exists exactly when
\begin{equation}
 \mathsf b(x_1)=\mathsf b(x_2)\ \Longrightarrow\ m(x_1)=m(x_2).
 \label{eq:reading-fibre}
\end{equation}
Indeed, a decoder $d$ with $m=d\circ\mathsf b$ implies the condition;
conversely define $d(\mathsf b(x))=m(x)$. The condition makes this
independent of the representative. Only the chosen label, not the entire
history or state, must be recovered.

For a fixed input $\iota$, let $T_\iota:X\to X$ be a state update. A
closed record update $U_\iota$ exists precisely when
\begin{equation}
 \mathsf b(x_1)=\mathsf b(x_2)\ \Longrightarrow\\
 \mathsf b(T_\iota x_1)=\mathsf b(T_\iota x_2).
 \label{eq:closure-fibre}
\end{equation}
The proof is the same: set $U_\iota(\mathsf b(x))=\mathsf b(T_\iota x)$.
For an autonomous semigroup $T_t$, if this holds for every $t$, the induced
updates form a semigroup because
$U_s U_t\mathsf b=\mathsf b T_s T_t=U_{s+t}\mathsf b$ and $\mathsf b$
is onto its image. No continuity or numerical stability follows merely
from this set-theoretic argument. It is the elementary factorization
underlying the more structured smooth reduction studied in
Ref.~\cite[Section 2.2]{horstmeyer}.

\section{Equal-energy labels and a finite reading window}
\label{sec:reading}

\subsection{A surface observable and the transfer matrix}

Assume access to $q$ and its time rate, and define
\begin{equation}
 z=\begin{pmatrix}\Omega q\\ \dot q+\gamma q/2\end{pmatrix},
 \qquad Y=\norm{z}^2=(\dot q+\gamma q/2)^2+\Omega^2q^2.
 \label{eq:readout}
\end{equation}
Both quantities use no instantaneous bulk trace. The scalar is invariant
under a simultaneous sign change of the real field and its rate;
Section~\ref{sec:reference} supplies its gauge-covariant description.
Access to the two ingredients of $Y$ is explicitly broader than access
to $q^2$ alone. A nonzero underdamped free solution has recurring zeros
of $q$, whereas $Y$ below has no such blind times.

Equation~\eqref{eq:oscillator} becomes
\begin{equation}
 \dot z=Lz+kF(t),\qquad
 L=\begin{pmatrix}-\gamma/2&\Omega\\-\Omega&-\gamma/2\end{pmatrix},
 \qquad k=\begin{pmatrix}0\\2\omega_*^2\end{pmatrix}.
 \label{eq:state}
\end{equation}
Direct differentiation verifies
\begin{align}
 \Phi(s)&=e^{-\gamma s/2}
 \begin{pmatrix}\cos(\Omega s)&\sin(\Omega s)\\
 -\sin(\Omega s)&\cos(\Omega s)\end{pmatrix},\\
 \Phi(s+t)&=\Phi(s)\Phi(t),\qquad
 \Phi(s)^{\mathsf T}\Phi(s)=e^{-\gamma s}I.
 \label{eq:propagator}
\end{align}
Consequently $Y(t_p+s)=e^{-\gamma s}Y(t_p)$ whenever no incident
input is present after $t_p$. More generally,
\begin{equation}
 z(t)=\Phi(t-t_0)z(t_0)+\int_{t_0}^t\Phi(t-s)kF(s)\,ds.
 \label{eq:state-convolution}
\end{equation}
This is the usual linear state-space solution\cite{frazzoli}, here
verified by differentiation and the initial value at $t_0$.

Fix a base pulse $f$ supported in $[0,T]$ and write
\begin{equation}
 c=\int_0^T\Phi(T-s)k f(s)\,ds,\qquad C_* =\norm{c}^2.
 \label{eq:unit-response}
\end{equation}
We require $C_*>0$. Nonzero input alone would not establish this:
the terminal integral could cancel for some waveforms. The explicit
pulse used below has a proved positive lower bound.

\subsection{Two orders with equal initial energy}

Choose $0<a<b$ and $\Delta\ge T$. The inputs are
\begin{align}
 F_{ab}(t)&=a f(t)+b f(t-\Delta),\\
 F_{ba}(t)&=b f(t)+a f(t-\Delta),\qquad t_p=T+\Delta.
 \label{eq:orders}
\end{align}
Their differentiated supports are disjoint except at vanishing endpoints.
Proposition~\ref{prop:energy} therefore gives
\begin{equation}
 E_{ab}=E_{ba}=(a^2+b^2)E_f,\qquad
 E_f=K\int_0^T f'(s)^2\,ds.
 \label{eq:equal-energy}
\end{equation}
Equal energy does not mean equal preparation. The spatial distribution
of the initial bulk pulses encodes their order with respect to the supplied
clock. For this pair even the modulus profiles on the separated supports
differ, so a phase transformation does not identify the preparations.

\begin{proposition}[Equal-energy order readout]
\label{prop:order}
Let $f$ be an admissible pulse with $C_*>0$, let $0<a<b$ and
$\Delta\ge T$, and assume $\Omega^2>0$. For the two ideal preparations
\eqref{eq:orders}, with no further incoming pulse after $t_p=T+\Delta$,
every $t\ge t_p$ satisfies
\begin{equation}
 Y_{ab}(t)-Y_{ba}(t)
 =(b^2-a^2)(1-e^{-\gamma\Delta})C_*
      e^{-\gamma(t-t_p)}>0.
 \label{eq:gap}
\end{equation}
\end{proposition}
\begin{proof}
At $t_p$, the first pulse response has propagated freely for $\Delta$
more units of time than the second, so
\[
 z_{ab}=a\Phi(\Delta)c+bc,\qquad
 z_{ba}=b\Phi(\Delta)c+ac.
\]
Set $\rho=e^{-\gamma\Delta/2}$ and $\vartheta=\Omega\Delta$.
Since $c^{\mathsf T}\Phi(\Delta)c=\rho\cos\vartheta\,C_*$,
\begin{align*}
 Y_{ab}(t_p)&=C_*(a^2\rho^2+b^2+2ab\rho\cos\vartheta),\\
 Y_{ba}(t_p)&=C_*(b^2\rho^2+a^2+2ab\rho\cos\vartheta).
\end{align*}
The cross terms cancel, leaving $(b^2-a^2)(1-\rho^2)C_*$.
Free propagation multiplies both squared norms by
$e^{-\gamma(t-t_p)}$, proving \eqref{eq:gap}.
\end{proof}

The algebra identifies the mechanism: the earlier contribution has aged
more before the common reading time. A quadratic scalar readout retains
this difference even when the total input energies coincide. No
nonlinearity of the evolution is needed.

\subsection{A bounded-error decoder}

At a known time $t$, suppose the measured value is
$\widetilde Y=Y_j(t)+\eta$, where $j\in\{ab,ba\}$ and $|\eta|\le\delta$.
The decoder knows the parameters, pulse shape, amplitudes and time.
Calibration uncertainty, if present, must be included in the declared
scalar bound; it is not silently set to zero after being incurred.

\begin{lemma}[Separation of reading intervals]
\label{lem:decoder}
A decoder correct for both labels and every allowed scalar error exists
if and only if $|Y_{ab}(t)-Y_{ba}(t)|>2\delta$. If the two calibrated
signals are continuous on a compact interval, correctness at every
specified time in that interval is equivalent to their minimum
separation being strictly greater than $2\delta$.
\end{lemma}
\begin{proof}
The allowed readings are the closed intervals
$[Y_j(t)-\delta,Y_j(t)+\delta]$. They are disjoint exactly under the
strict inequality; their touching endpoint at equality is already
ambiguous. If disjoint, comparing with the midpoint of the centres
decodes the two labels. Continuity on a compact interval makes the
minimum attainable, proving the uniform assertion.
\end{proof}

With $D_0=(b^2-a^2)(1-e^{-\gamma\Delta})C_*$, the decreasing gap makes
its value at the last reading time the relevant one. For $\delta>0$
and $D_0>2\delta$, the strict deadline is
\begin{equation}
 t<t_p+\frac1\gamma\log\frac{D_0}{2\delta}.
 \label{eq:deadline}
\end{equation}
If $D_0\le2\delta$, no post-pulse time meets this error criterion.
When $\delta=0$, the labels remain distinct at every finite time, but
their asymptotic gap is still zero. A noise-dependent decoding deadline
and an asymptotically nonzero residual are different claims.

\subsection{One explicit analytic margin}
\label{sec:margin}

Choose dimensionless parameters $K=M=g=1$ and $T=1$. This selects a
coupling ratio $gM/K^2=1$; it is not merely a change of units removing
all possible dimensionless parameters. Let
\begin{equation}
 f(t)=256t^4(1-t)^4\quad(0\le t\le1),\qquad f(t)=0\quad\hbox{otherwise}.
 \label{eq:pulse}
\end{equation}
The zero extension is $C^3$. The derivative is
$f'(t)=1024t^3(1-t)^3(1-2t)$. Expanding $(1-2t)^2$ and integrating
monomials, or repeatedly integrating by parts, gives
\begin{align}
 E_f&=2^{20}\int_0^1t^6(1-t)^6(1-2t)^2\,dt
       =\frac{262144}{45045},\\
 \int_0^1(1-s)f(s)\,ds&=256\frac{4!\,5!}{10!}=\frac{64}{315}.
 \label{eq:pulse-integrals}
\end{align}
Here the elementary identity
$\int_0^1t^m(1-t)^n\,dt=m!n!/(m+n+1)!$ for nonnegative integers
follows by integration by parts in $n$ followed by integrating $t^{m+n}$.

For $0\le z\le\sqrt3/2$, integration of $\cos z\ge1-z^2/2$ gives
$\sin z\ge z-z^3/6\ge7z/8$. Hence, for $0\le r\le1$,
\[
 G(r)=e^{-r/2}\frac{\sin((\sqrt3/2)r)}{\sqrt3/2}
       \ge\frac78e^{-1/2}r.
\]
The pulse is nonnegative, so the convolution yields
\begin{equation}
 q(1)\ge2\frac78e^{-1/2}\frac{64}{315}
       =\frac{16}{45\sqrt e},\qquad
 C_*\ge\frac34q(1)^2\ge\frac{64}{675e}>0.
 \label{eq:c-lower}
\end{equation}
This supplies the required nonzero terminal response without numerical
quadrature.

Now set $a=1$, $b=2$, $\Delta=2$, and declare the reading interval
$[3,7]$ with $\delta=10^{-4}$. Both initial energies are
$5E_f=262144/9009$. For every time in this interval,
\begin{align}
 Y_{12}(t)-Y_{21}(t)
 &\ge\frac{64}{225e^5}(1-e^{-2})\nonumber\\
 &>\frac{16}{75\,3^5}>\frac{2}{10^4}=2\delta.
 \label{eq:margin}
\end{align}
We used $2<e<3$, which follows from the factorial series, and hence
$1-e^{-2}>3/4$. The last comparison is rational. The midpoint decoder
therefore selects the order from a single reading at any specified time
in $[3,7]$.

This is a guarantee for ideal prepared data under a bounded scalar
observation error. It does not certify a detector's access to $\dot q$,
clock accuracy, pulse-shape accuracy or preparation errors. The statement
at any specified time can be tested on separate prepared repetitions;
no nondisturbing sequence of physical measurements is presumed. The
energy comparison concerns the mathematical field data, not an omitted
measurement apparatus.

\section{Why the scalar does not close, and what does}
\label{sec:closure}

\subsection{The next pulse exposes discarded orientation}

Keep the same base pulse with $C_*>0$. Between successive pulse end times
separated by $\Delta\ge T$, assume no incident input except one pulse
$\beta f$ occupying the final $T$ units. Then
\begin{equation}
 z_{n+1}=\Phi(\Delta)z_n+\beta c.
 \label{eq:affine}
\end{equation}
Indeed, the free contribution is the first term of
\eqref{eq:state-convolution}; a change of variables in its integral gives
exactly $\beta c$. Squaring gives
\begin{equation}
 Y_{n+1}=e^{-\gamma\Delta}Y_n+\beta^2C_*
             +2\beta c^{\mathsf T}\Phi(\Delta)z_n.
 \label{eq:scalar-update}
\end{equation}
The last term depends on the orientation of the current field--rate
vector relative to the incoming pulse, not only its squared norm.

To decide whether this dependence is genuinely an obstruction, we must
specify a family containing different states with the same scalar.
We now enlarge the preparation family: any finite sequence of separated
copies of $f$ with arbitrary real amplitudes is allowed. All are prepared
as initial incoming bulk profiles with the boundary initially at rest.
The signed amplitudes are relative to one common preparation standard.
The earlier two positive orders belong to this family, but do not exhaust it.

\begin{lemma}[Reachable terminal vectors]
\label{lem:reachable}
Every real vector in $\R^2$ is a terminal state of two separated signed
pulses of the fixed shape, at a common terminal time chosen as below.
\end{lemma}
\begin{proof}
Choose $s\ge T$ with $\sin(\Omega s)\ne0$. The two pulse supports can
be $[0,T]$ and $[s,s+T]$, and their terminal vector is
$\alpha\Phi(s)c+\beta c$. Since
\begin{equation}
 \det[c,\Phi(s)c]=-e^{-\gamma s/2}\sin(\Omega s)C_*\ne0,
 \label{eq:reachability}
\end{equation}
these two vectors form a basis. Solving the real two-by-two linear
system gives finite amplitudes for every desired vector. The preparation
has finite energy $(\alpha^2+\beta^2)E_f$ and satisfies the same initial
and endpoint compatibility conditions as before.
\end{proof}

\begin{proposition}[No universal scalar update on the enlarged family]
\label{prop:nonclosure}
Fix $\Delta\ge T$ and a next amplitude $\beta>0$. There is no function
$\mathcal U_\beta:[0,\infty)\to[0,\infty)$ such that
$Y_{n+1}=\mathcal U_\beta(Y_n)$ for every reachable terminal state in
Lemma~\ref{lem:reachable} followed by that same next pulse.
\end{proposition}
\begin{proof}
Let $d=\Phi(\Delta)^{\mathsf T}c\ne0$ and
$J=\left(\begin{smallmatrix}0&1\\-1&0\end{smallmatrix}\right)$.
For $R>0$ choose
\begin{equation}
 z^{(1)}=R\frac d{\norm d},\qquad
 z^{(2)}=R\frac{Jd}{\norm d}.
 \label{eq:witnesses}
\end{equation}
Their squared norms are both $R^2$, whereas
$d^{\mathsf T}z^{(1)}=R\norm d$ and $d^{\mathsf T}z^{(2)}=0$.
The next scalar values from \eqref{eq:scalar-update} differ by
$2\beta R\norm d>0$. Both states are realizable by
Lemma~\ref{lem:reachable}, and appending the same future pulse gives
admissible complete incoming profiles. Thus the fibre condition
\eqref{eq:closure-fibre} fails.
\end{proof}

The witnesses differ in orientation, not merely by a common sign.
Their complete initial energies have not been constrained to agree.
Equal energy was essential to the comparison in
Proposition~\ref{prop:order}, but is not an assumption of this
nonclosure proposition. Adding that constraint would require a separate
argument.

In particular, the counterexample does not invalidate reading the two
positive order labels. At a fixed reading time those labels have
\emph{different} scalars. On a family where a scalar uniquely identifies
the allowed current state, one can predict the next scalar by identifying
that state first. The obstruction concerns the declared larger family;
it is not a prohibition of every finite-dimensional record.

\subsection{A common phase sector and a supplied reference}
\label{sec:reference}

The real vector $z$ is already sufficient for \eqref{eq:affine} in its
chosen real representative. To express that sufficiency without treating
an absolute charged phase or sign as an observable, specify the following
restricted complex description. On a trivial bundle let $\theta$ be
smooth and write
\begin{equation}
 \eta=e^{i\theta},\qquad \Psi=\eta u,\qquad \Xi=\eta q,
 \qquad A_\mu=\partial_\mu\theta,\qquad D_\mu=\partial_\mu-iA_\mu.
 \label{eq:common-phase}
\end{equation}
Then $D_\mu\Psi=\eta\partial_\mu u$ and $D_t\Xi=\eta\dot q$.
The curvature is zero, and currents proportional to
$\operatorname{Im}(\overline\Psi D_\mu\Psi)$ vanish. This is the common
phase orbit of the real neutral solutions, not arbitrary charged or
independently phased input data. Substitution into the covariant versions
of \eqref{eq:system} gives exactly their real counterparts.

Under a gauge transformation with phase $\chi$,
\[
 (\Psi,\Xi,\eta)\mapsto e^{i\chi}(\Psi,\Xi,\eta),\qquad
 A_\mu\mapsto A_\mu+\partial_\mu\chi.
\]
The surface reference $\eta$ is supplied by preparation and calibration,
with $|\eta|=1$ and $D_t\eta=0$. It is not constructed by normalizing
$\Xi$, and remains defined when $q=0$. All future signed amplitudes are
prepared relative to this same reference. Its physical maintenance and
measurement are not supplied by a reference-field action here.

Define
\begin{equation}
 w=\begin{pmatrix}\Omega\Xi\\D_t\Xi+\gamma\Xi/2\end{pmatrix},
 \qquad r=\overline\eta\,w\in\R^2.
 \label{eq:relative}
\end{equation}
Both components of $w$ acquire the phase $e^{i\chi}$; its cancellation
against the transformed $\overline\eta$ makes $r$ gauge invariant.
The corresponding scalar is
\begin{equation}
 Y=w^\dagger w=\norm r^2
   =|D_t\Xi+\gamma\Xi/2|^2+\Omega^2|\Xi|^2,
 \label{eq:gauge-scalar}
\end{equation}
which reduces to \eqref{eq:readout}. The scalar itself does not require
reading the external phase standard. The standard is relevant to the
larger record and its relation to later inputs.

\begin{proposition}[Orbit separation for the reduced record]
\label{prop:orbits}
On
\[
 \mathcal P=\{(w,\eta): |\eta|=1,\ \overline\eta w\in\R^2\},
\]
$r$ distinguishes the orbits of the simultaneous phase action on $(w,\eta)$.
It need not distinguish full bulk states or complete input histories.
\end{proposition}
\begin{proof}
Invariance has just been shown. If two pairs have the same $r$, then
$w=\eta r$ and $w'=\eta'r$. The unit factor $\eta'/\eta$ maps both
members of the first pair to the second. Conversely a simultaneous phase
change leaves $r$ unchanged. The proof includes $w=0$ because $\eta$
never vanishes.
\end{proof}

This quotient statement neither derives access to $r$ from access to $Y$
nor proves that a phase reference is energetically free. It identifies
the additional relation that is being made available.

\subsection{Controlled closure and its limits}

Along the surface between successive pulse end times, let
\begin{equation}
 P_{n+1,n}=\exp\left(i\int_{t_n}^{t_{n+1}}A_0(t)\,dt\right),
 \qquad \eta_{n+1}=P_{n+1,n}\eta_n.
 \label{eq:transport}
\end{equation}
The positive exponent follows by integrating $D_t\eta=0$.
Under a gauge transformation the transport acquires
$e^{i\chi_{n+1}}$ on the left and $e^{-i\chi_n}$ on the right.
For the same input convention as in \eqref{eq:affine}, the covariant
update is
\begin{equation}
 w_{n+1}=P_{n+1,n}\Phi(\Delta)w_n+
                     \eta_{n+1}\beta_n c.
 \label{eq:covariant-update}
\end{equation}
It follows by multiplying the real update by $\eta_{n+1}$ and using
$w_n=\eta_n z_n$. It is established only on the stipulated common-phase
family.

\begin{proposition}[Closed update with the reference and input supplied]
\label{prop:relative-closure}
On the common-phase family \eqref{eq:common-phase}, supply the parallel
reference \eqref{eq:transport}. Between pulse end times separated by
$\Delta\ge T$, admit only the prescribed next pulse $\beta_n f$ during
the final $T$ units, with no other incoming excitation. Then the relative
record has the exact update
\begin{equation}
 r_{n+1}=\Phi(\Delta)r_n+\beta_n c.
 \label{eq:relative-update}
\end{equation}
Thus equal relative records with the same next pulse protocol give equal
next relative records, and $Y_{n+1}=\norm{r_{n+1}}^2$ is determined.
\end{proposition}
\begin{proof}
Multiply \eqref{eq:covariant-update} by $\overline\eta_{n+1}$.
Since $|P_{n+1,n}|=1$ and
$\overline\eta_{n+1}P_{n+1,n}=\overline\eta_n$, the transport cancels
and \eqref{eq:relative-update} follows.
\end{proof}

Even the intrinsic matrix $w w^\dagger$ would not supply every relation
to an independently specified next input. With the reference held fixed,
$r$ and $-r$ give the same matrix but their next squared norms differ by
$4\beta c^{\mathsf T}\Phi(\Delta)r$, which is nonzero, for example, at
$r=\Phi(\Delta)^{\mathsf T}c$. Such states are reachable by the signed
preparations. Negating the past alone is not the simultaneous phase
change of the past, reference and future input that gauge invariance
would identify.

The enlarged record is sufficient, not proved minimal among all possible
access rules. Nor is its controlled closure autonomous from the bulk:
$\beta_n$, timing, shape and reference relation are specified resources.
For example, $r_n=0$ followed by no pulse gives a different next state
from $r_n=0$ followed by $\beta_n\ne0$. Incorporating a future schedule
in an enlarged state records that schedule; it does not generate it.

A common pulse also cannot restore a lost vector separation in this
linear model. For two histories receiving the same next input,
\begin{equation}
 r_{n+1}^{(1)}-r_{n+1}^{(2)}
  =\Phi(\Delta)(r_n^{(1)}-r_n^{(2)}),\qquad
 \norm{r_{n+1}^{(1)}-r_{n+1}^{(2)}}
  =e^{-\gamma\Delta/2}\norm{r_n^{(1)}-r_n^{(2)}}.
 \label{eq:contraction}
\end{equation}
The common affine contribution cancels. A quadratic scalar gap may
change because of its cross term with that contribution, but this is
not amplification of the full vector record.

\section{Interpretation, related work and limits}
\label{sec:discussion}

\subsection{The scope of the three statements}

Table~\ref{tab:scope} collects the domains rather than combining the
results into one claim with unstated hypotheses. The fixed dynamics
are shared, but the allowed preparations and available information
are not identical in the three rows.

\begin{table}[htbp]
\centering
\small
\renewcommand{\arraystretch}{1.2}
\begin{tabularx}{\textwidth}{@{}>{\raggedright\arraybackslash}p{0.19\textwidth}>{\raggedright\arraybackslash}X>{\raggedright\arraybackslash}X@{}}
\toprule
Statement & Family and supplied access & What is not implied \\
\midrule
Order reading & Two positive orders at equal total energy; known time,
calibration and scalar $Y$ using field and rate & Decoding arbitrary
histories or scalar closure on a larger family \\
Scalar nonclosure & Reachable signed-pulse histories; a common next
input; the same current scalar & Nonclosure on the two-label family,
equal energies of the witnesses, or failure of every finite record \\
Relative closure & Common-phase family; supplied reference and next
pulse protocol; two real components & Minimality, autonomous input
generation or persistent preservation of old labels \\
\bottomrule
\end{tabularx}
\caption{Preparation and access assumptions are part of each statement.}
\label{tab:scope}
\end{table}

The role of the field theory is concrete but limited. It supplies
compatible finite-energy histories and a conservative account of the
outgoing radiation. The observable distinction and the obstruction to
scalar updating are then read from a two-dimensional linear system.
One should not replace that explanation by the assertion that an
independent physical memory mechanism has been discovered.

\subsection{Comparison with established constructions}

A direct model comparison is available with George's boundary oscillator
system\cite[Section 2.1, Eqs.~(14), (19)--(21)]{george}. In canonically
normalized coordinates $\varphi=\sqrt K\,u$ and $Q=\sqrt M\,q$, our
Hamiltonian has the same form under the parameter specialization
\begin{equation}
 m=0,\qquad \lambda=\frac gK,\qquad
 \mu^2=\frac gM,\qquad \beta_{\rm G}=-\frac g{\sqrt{KM}}.
 \label{eq:george-map}
\end{equation}
Here $\beta_{\rm G}$ denotes that paper's coupling, not a pulse amplitude.
Expanding our interaction energy verifies the comparison directly.
The referenced treatment assumes positive bulk mass; this is a massless
specialization of its displayed equations, not an import of its full
spectral or quantum conclusions. The parameter map satisfies
$\beta_{\rm G}^2=\lambda\mu^2$, so at $m=0$ George's bound-state
polynomial factors as $\rho(\rho^2+\lambda\rho+\mu^2)$
\cite[Section~2.2]{george}. Its zero root
is not a positive inverse decay length and hence is not a localized
bound-mode root. Our pulse solution and energy balance were proved
independently. This comparison locates the model in an existing class
rather than claiming a new coupling mechanism.

Earlier brane--bulk work also studies one-dimensional toy models coupling
field dynamics to discrete boundary degrees of freedom, including
string--oscillator analogues\cite{binetruy}. This is a close conceptual
precedent, though not the same independent boundary coordinate and spring
coupling used here.

Related field--mechanical systems have also been studied in interval
models with endpoint point masses; Barbero et al. analyze their modes
and quantization\cite{barbero}. Dynamical Wentzell conditions for a
Klein--Gordon field in anti-de Sitter spacetime provide another established
bulk--boundary setting\cite{dappiaggi}. These works differ in geometry
and boundary variables from the present massless half-line model. They
confirm that boundary coupling and mode analysis are not the new result
claimed here. The narrower question in this paper is the explicit
pulse-order readout and scalar-update counterexample under the stated
preparation and access rules; no priority claim is made for those
constructions.

The state-space convolution and affine update are standard linear
operations\cite{frazzoli}. In the present example the order effect is
obtained by cancelling the common cross terms in two squared norms.
The explicit energy comparison and analytic error margin clarify the
protocol, but should not be presented as a new law of linear response.

Horstmeyer and Atay characterize exact lumpability of smooth vector
fields\cite{horstmeyer}. Their Proposition~2 relates the original and
reduced flows by a commuting projection, with attention to the domains
of definition. Equation~\eqref{eq:closure-fibre} gives the elementary
representative-independence condition needed here. It does not include
their smoothness structure and is not a replacement theorem for that
setting. The contribution of Proposition~\ref{prop:nonclosure} is the
explicit reachable pair that violates the condition for this observable.

Shalizi and Crutchfield construct causal states for prediction of
stochastic processes and establish optimality and uniqueness
properties\cite{shalizi}. Barnett and Crutchfield extend that approach
to input--output processes and channel models\cite[Sections VI and IX]{barnett}.
These provide the relevant distinction between a summary of the past
and information sufficient for future prediction under specified inputs.
Our deterministic finite pulse families have neither a probability law
for input histories nor an information-theoretic minimization criterion.
Consequently the relative vector is not identified with a proved minimal
causal-state representation. The analogy motivates a careful access
statement; it does not transfer those optimality results.

\subsection{Resources, limitations and conclusion}

The model, boundary location, clock, parameters and allowed input family
are supplied rather than derived from Field 01. The decoder additionally
uses calibrated signals at known times and assumes access to a
rate-dependent scalar. The enlarged record requires a phase reference
related to the incident preparation. The analysis does not test the
cosmological claims of Field 01. The word \emph{record} is used
operationally for retained data supporting a declared reading or update
task, not as evidence of such a principle.

The explicit guarantee is finite-time and has a declared deterministic
error model. It does not establish robustness to arbitrary charged
perturbations, finite-bandwidth detection, unknown timing or unknown
waveforms. Even the ideal signal gap decays. The fully conservative
bulk--boundary system can therefore support the stated temporary order
code without providing permanent surface storage.

The construction isolates one useful distinction. A scalar surface
observation reads a chosen equal-energy order label, but loses orientation
needed for a common future input when the preparation family is enlarged.
A relative two-component record retains that orientation and closes the
specified controlled update. The difference between these statements is
not a numerical precision issue: it is the difference between a label
map on a restricted family and a state description adequate for an
announced class of future operations.

\begin{thebibliography}{9}
\bibitem{george}
A.~George, ``The Massive Klein--Gordon Field Coupled to a Harmonic
Oscillator at the Boundary,'' \emph{Journal of Physics A: Mathematical
and General} \textbf{38} (2005), 7399--7418.
\url{https://doi.org/10.1088/0305-4470/38/33/013}.
Comparison uses Sections~2.1--2.2 of \texttt{hep-th/0412067v3},
26 September 2005.

\bibitem{binetruy}
P.~Binetruy, M.~Bucher and C.~Carvalho, ``Models of the Brane--Bulk
Interaction: Toward Understanding Braneworld Cosmological Perturbations,''
\emph{Physical Review D} \textbf{70}, 043509 (2004).
\url{https://arxiv.org/abs/hep-th/0403154}.

\bibitem{barbero}
J.~F.~Barbero G., B.~A.~Ju\'arez-Aubry, J.~Margalef-Bentabol and
E.~J.~S.~Villase\~nor, ``Quantization of scalar fields coupled to
point-masses,'' preprint, \url{https://arxiv.org/abs/1501.05114}.

\bibitem{dappiaggi}
C.~Dappiaggi, H.~R.~C.~Ferreira and B.~A.~Ju\'arez-Aubry,
``Mode solutions for a Klein--Gordon field in anti-de Sitter spacetime
with dynamical boundary conditions of Wentzell type,''
\emph{Physical Review D} \textbf{97}, 085022 (2018).
\url{https://doi.org/10.1103/PhysRevD.97.085022}.
Preprint: \url{https://arxiv.org/abs/1802.00283}.

\bibitem{frazzoli}
E.~Frazzoli, \emph{Solutions of State-space Models}, Lecture~8,
6.241J Dynamic Systems and Control, Massachusetts Institute of Technology,
28 February 2011.
\href{https://ocw.mit.edu/courses/6-241j-dynamic-systems-and-control-spring-2011/resources/mit6_241js11_lec08/}{MIT OpenCourseWare, lecture notes}.

\bibitem{horstmeyer}
L.~Horstmeyer and F.~M.~Atay, ``Characterization of Exact Lumpability
for Vector Fields on Smooth Manifolds,''
\emph{Differential Geometry and its Applications} \textbf{48} (2016), 46--60.
\url{https://doi.org/10.1016/j.difgeo.2016.06.001}.
See also \texttt{arXiv:1607.01237v2}, Section~2.2, Proposition~2.

\bibitem{shalizi}
C.~R.~Shalizi and J.~P.~Crutchfield, ``Computational Mechanics:
Pattern and Prediction, Structure and Simplicity,''
\emph{Journal of Statistical Physics} \textbf{104} (2001), 817--879.
\url{https://doi.org/10.1023/A:1010388907793}.

\bibitem{barnett}
N.~Barnett and J.~P.~Crutchfield, ``Computational Mechanics of
Input-Output Processes: Structured Transformations and the
$\epsilon$-Transducer,'' \texttt{arXiv:1412.2690v3}, 26 January 2016.
\url{https://arxiv.org/abs/1412.2690v3}.
Published article: \url{https://doi.org/10.1007/s10955-015-1327-5}.
Sections~VI and IX concern the transducer and its optimality.
\end{thebibliography}
\end{document}